\documentclass[a4paper,14pt]{report}
\usepackage{graphicx}
\usepackage[T2A]{fontenc}
\usepackage[utf8]{inputenc}
\usepackage[russian]{babel}
\usepackage{amsmath, amssymb, amsthm}
\usepackage{hyperref}
\usepackage{geometry}
\geometry{top=2cm, bottom=2cm, left=2.5cm, right=1.5cm}
\usepackage{listings}
\usepackage[table]{xcolor}
\usepackage{setspace}
\onehalfspacing
\setcounter{secnumdepth}{-3}

\title{Адаптационный курс 2026}

\begin{document}

\maketitle
\tableofcontents

\part{Глава №1. Математические основы функционирования ЭВМ.}

\chapter{Тема №1. Системы счисления.}
\section{Лекция №1. Позиционные и непозиционные СС. Перевод из 10-й в N-ю и обратно. Быстрый перевод из 2-й в 4-, 8-, 16-ю. Базовая целочисленная арифметика в N-й СС.}

\subsection{Что такое система счисления}
\textbf{Система счисления (СС)} --- совокупность символов (цифр) и правил для наименования и обозначения чисел.

\textbf{Алфавит системы счисления} --- совокупность цифр.
\textbf{Мощность алфавита} --- количество цифр в алфавите.

Системы счисления подразделяются на позиционные и непозиционные. В непозиционной Римской системе счисления значение цифры не зависит от её позиции в числе: I -- 1, II -- 2, X -- 10.

В позиционных системах счисления числовое значение цифры определяется её позицией в числе. Все позиционные системы счисления имеют основание, равное мощности алфавита:

\begin{itemize}
    \item двоичная (СС$_{2}$): 0, 1
    \item восьмеричная (СС$_{8}$): 0, 1, 2, 3, 4, 5, 6, 7
    \item десятичная (СС$_{10}$): 0, 1, 2, 3, 4, 5, 6, 7, 8, 9
    \item шестнадцатеричная (СС$_{16}$): 0, 1, 2, 3, 4, 5, 6, 7, 8, 9, A, B, C, D, E, F.
\end{itemize}

В общем виде число в позиционной системе счисления записывается в виде полинома:
\[
A = \sum_{i=-m}^{n-1} a_i \cdot S^i,
\]
где $S$ --- основание системы, $a_i$ --- одна из цифр алфавита этой системы счисления.

При сокращённой записи числа основание подразумевается, а цифры целой и дробной части разделяются точкой.

Каждый левый разряд имеет вес в $S$ раз больший, чем предыдущий. Крайний правый разряд называется младшим, крайний левый --- старшим.

В памяти ЭВМ числа хранятся в пределах разрядной сетки с фиксированным числом разрядов. Если на запись числа без знака отведено $n$ разрядов на целую часть и $m$ на дробную, то максимальное и минимальное числа:
\[
A_{\max} = (S^n - 1) + (1 - S^{-m}), \qquad A_{\min} = S^{-m}.
\]
Количество различных чисел в такой сетке: $N = S^{m+n}$.

\subsection{Преобразования из одних систем счисления в другие}
В основе любого преобразования лежит представление числа в виде полинома.

\textbf{Перевод из любой СС в десятичную}: вычислить значение полинома, подставляя десятичные эквиваленты цифр и основания.

Примеры:
\begin{align*}
101.11_2 &= 1\cdot2^2 + 0\cdot2^1 + 1\cdot2^0 + 1\cdot2^{-1} + 1\cdot2^{-2} = 4+0+1+0.5+0.25 = 5.75_{10},\\
123.4_8 &= 1\cdot8^2 + 2\cdot8^1 + 3\cdot8^0 + 4\cdot8^{-1} = 64+16+3+0.5 = 83.5_{10},\\
AB.8_{16} &= 10\cdot16^1 + 11\cdot16^0 + 8\cdot16^{-1} = 160+11+0.5 = 171.5_{10}.
\end{align*}

\textbf{Перевод из десятичной в любую СС}:
\begin{itemize}
    \item \textit{Целая часть:} последовательно делить на основание новой СС, записывая остатки; последний результат и остатки в обратном порядке дают число.
    \item \textit{Дробная часть:} последовательно умножать дробную часть на основание, записывая целые части в прямом порядке; процесс останавливается при нулевой дробной части или достижении нужной точности.
\end{itemize}

Пример перевода $23_{10}$ в двоичную:
\[
23/2 = 11\ (1),\ 11/2=5\ (1),\ 5/2=2\ (1),\ 2/2=1\ (0),\ 1/2=0\ (1) \ \Rightarrow\ 10111_2.
\]
Перевод $0.25_{10}$ в двоичную:
\[
0.25\cdot2=0.5\ (0),\ 0.5\cdot2=1.0\ (1) \ \Rightarrow\ 0.01_2.
\]

\subsection{Быстрый перевод между двоичной, восьмеричной и шестнадцатеричной системами}
Поскольку $8=2^3$ и $16=2^4$, перевод сводится к группировке двоичных разрядов.

\begin{itemize}
    \item \textbf{Из двоичной в восьмеричную:} разбить число на триады (группы по 3 бита) от запятой влево и вправо, заменить каждую триаду восьмеричной цифрой.
    \item \textbf{Из восьмеричной в двоичную:} каждую восьмеричную цифру заменить трёхбитовой двоичной триадой.
    \item \textbf{Из двоичной в шестнадцатеричную:} разбить на тетрады (по 4 бита), заменить каждую тетраду шестнадцатеричной цифрой.
    \item \textbf{Из шестнадцатеричной в двоичную:} каждую шестнадцатеричную цифру заменить четырёхбитовой тетрадой.
\end{itemize}

Пример: $1111110.11001_2 \to$ восьмеричная: \texttt{001 111 110 . 110 010} $\to 176.62_8$.
В шестнадцатеричную: \texttt{0111 1110 . 1100 1000}? (дополняем нулями) $\to 7E.C8_{16}$.

Перевод между восьмеричной и шестнадцатеричной обычно выполняют через двоичную.

\subsection{Базовая целочисленная арифметика в N-й СС}
Арифметические операции во всех позиционных системах счисления основаны на таблицах сложения, вычитания и умножения одноразрядных чисел. Рассмотрим двоичную арифметику:

\textit{Сложение:}
\begin{center}
\begin{tabular}{c|c|c|c}
0+0=0 & 0+1=1 & 1+0=1 & 1+1=10
\end{tabular}
\end{center}

\textit{Вычитание:}
\begin{center}
\begin{tabular}{c|c|c|c}
0-0=0 & 10-1=1 & 1-1=0 & 1-0=1
\end{tabular}
\end{center}

\textit{Умножение:}
\begin{center}
\begin{tabular}{c|c|c|c}
0$\cdot$0=0 & 0$\cdot$1=0 & 1$\cdot$0=0 & 1$\cdot$1=1
\end{tabular}
\end{center}

Пример сложения: $11101.01_2 + 1011.10_2 = 101000.11_2$ (подробно поразрядно).

Пример вычитания: $11001_2 - 1011_2 = 1110_2$.

Умножение выполняется как последовательное умножение множимого на цифры множителя и сложение частичных произведений.

Деление — последовательное умножение цифр частного на делитель и вычитание.

В восьмеричной и шестнадцатеричной системах правила те же, но основание равно 8 или 16.

\section{Семинар №1. Перевод из 10-й в N-ю и обратно. Быстрый перевод из 2-й в 4-, 8-, 16-ю. Базовая целочисленная арифметика в N-й СС.}

\subsection*{Задача 1 (перевод целых чисел $10 \to N$)}
Переведите число $250_{10}$ в двоичную, восьмеричную и шестнадцатеричную системы.

\textbf{Решение:}
\begin{itemize}
    \item В двоичную: $250 / 2 = 125 (0)$, $125/2=62 (1)$, $62/2=31 (0)$, $31/2=15 (1)$, $15/2=7 (1)$, $7/2=3 (1)$, $3/2=1 (1)$, $1/2=0 (1)$ $\to 11111010_2$.
    \item В восьмеричную: $250 / 8 = 31 (ост. 2)$, $31/8=3 (ост. 7)$, $3/8=0 (ост. 3)$ $\to 372_8$.
    \item В шестнадцатеричную: $250 / 16 = 15 (ост. 10=A)$, $15/16=0 (ост. 15=F)$ $\to FA_{16}$.
\end{itemize}

\subsection*{Задача 2 (перевод $N \to 10$)}
Даны числа: $1100101_2$, $145_8$, $A9_{16}$. Переведите их в десятичную систему.

\textbf{Решение:}
\begin{align*}
1100101_2 &= 1\cdot64 + 1\cdot32 + 0\cdot16 + 0\cdot8 + 1\cdot4 + 0\cdot2 + 1\cdot1 = 101_{10},\\
145_8     &= 1\cdot64 + 4\cdot8 + 5\cdot1 = 101_{10},\\
A9_{16}   &= 10\cdot16 + 9 = 169_{10}.
\end{align*}

\subsection*{Задача 3 (быстрый перевод $2\leftrightarrow 8$, $2\leftrightarrow 16$)}
Дано двоичное число $110111001010_2$. Переведите его в восьмеричную и шестнадцатеричную системы.

\textbf{Решение:}
Группы по 3: \texttt{110 111 001 010} $\to$ \texttt{6 7 1 2} $\to 6712_8$.
Группы по 4: \texttt{1101 1100 1010} $\to$ \texttt{D C A} $\to DCA_{16}$.

Дано $2F5_{16}$. Переведите в двоичную и восьмеричную.

\textbf{Решение:}
$2=0010$, $F=1111$, $5=0101$ $\to 001011110101_2$ (убираем ведущие нули: $1011110101_2$).
В 8-ю: разбиваем на триады: \texttt{1 011 110 101} $\to$ \texttt{001 011 110 101} $\to$ \texttt{1 3 6 5} $\to 1365_8$.

\subsection*{Задача 4 (арифметика)}
Выполните действия в указанных системах. Проверьте переводом в десятичную.
\begin{enumerate}
    \item[a)] $1011_2 + 1101_2$
    \item[b)] $11001_2 - 1011_2$
    \item[c)] $23_8 \times 5_8$
    \item[d)] $1A3_{16} + 4F_{16}$
    \item[e)] $200_8 / 7_8$ (целочисленное деление)
\end{enumerate}

\textbf{Решение:}
\begin{enumerate}
    \item[a)] $1011_2 (11) + 1101_2 (13) = 11000_2 (24)$.
    \item[b)] $11001_2 (25) - 1011_2 (11) = 1110_2 (14)$.
    \item[c)] $23_8 \times 5_8 = 137_8$ (проверка: $19\cdot5=95=137_8$).
    \item[d)] $1A3_{16} (419) + 4F_{16} (79) = 498_{10} = 1F2_{16}$.
    \item[e)] $200_8 (128) / 7_8 (7) = 18$ (ост. 2) $\to$ частное $22_8$, остаток 2.
\end{enumerate}


\chapter{Тема №2. Представление чисел в памяти компьютера.}
\section{Лекция №2. Целые числа: знаковые и беззнаковые. Диапазоны. Прямой, обратный, дополнительный коды. Оптимизация арифметики отрицательных чисел с помощью доп. кода. Побитовые операции. Двоичная арифметика.}

\subsection{Представление беззнаковых целых чисел (unsigned)}
В ячейках памяти информация хранится в битах. Для беззнаковых целых все биты представляют значение. Если под число отведено $n$ бит, диапазон значений: $0 \dots 2^n-1$.

Пример: $n=8$ (байт). Диапазон $0\dots255$. Число $153_{10} = 10011001_2$.

Арифметические операции над беззнаковыми выполняются обычной двоичной арифметикой; переполнение фиксируется флагом переноса \textbf{CF}.

\subsection{Представление знаковых целых чисел}
Для представления отрицательных чисел используют старший бит как \textbf{знаковый}: 0 — плюс, 1 — минус. Существуют три основных кода: прямой, обратный, дополнительный.

\subsubsection{Прямой код}
Старший бит — знак, остальные — модуль числа.
Пример (8 бит): $+18 = 00010010_2$, $-18 = 10010010_2$.

Диапазон: $-(2^{n-1}-1) \dots +(2^{n-1}-1)$. Два нуля: $+0$ и $-0$. Арифметика требует анализа знаков, поэтому неудобен.

\subsubsection{Обратный код}
Отрицательное число получается побитовой инверсией модуля (все 0 заменяются 1, 1 на 0), включая знаковый бит.
Пример: $+18 = 00010010_2$; $-18 = 11101101_2$.

Диапазон такой же, как у прямого. Два нуля: $+0 = 00\dots0$, $-0 = 11\dots1$. Сложение требует циклического переноса.

\subsubsection{Дополнительный код}
Отрицательное число формируется путём инверсии модуля и прибавления 1 к младшему разряду (либо вычитанием модуля из $2^n$).
Пример: $+18 = 00010010_2$; инверсия $11101101_2$, $+1 \to 11101110_2 = -18$. Также можно: $2^8 - 18 = 238 = 11101110_2$.

Диапазон: $-2^{n-1} \dots 2^{n-1}-1$. Ноль единственный: $00\dots0$. Дополнительный код стал стандартом, потому что позволяет выполнять вычитание как сложение: $A - B = A + (-B)$.

\subsection{Оптимизация арифметики отрицательных чисел с помощью доп. кода}
АЛУ использует только сумматор. Пример: $5 - 3 = 5 + (-3)$: $00000101 + 11111101 = 00000010$ (перенос за пределы разрядной сетки игнорируется).

При расширении разрядности используется знаковое расширение (копирование знакового бита).

\subsection{Побитовые операции}
\begin{itemize}
    \item \textbf{NOT (побитовое отрицание)}: $\sim$ — инвертирует каждый бит.
    \item \textbf{AND (И)}: 1, если оба бита 1.
    \item \textbf{OR (ИЛИ)}: 1, если хотя бы один бит 1.
    \item \textbf{XOR (исключающее ИЛИ)}: 1, если биты различны.
    \item \textbf{Сдвиги:}
    \begin{itemize}
        \item \textit{Логический сдвиг влево} (<<): биты сдвигаются влево, справа нули. Умножение на $2^k$ (для беззнаковых).
        \item \textit{Логический сдвиг вправо} (>> для unsigned): сдвиг вправо, слева нули. Деление на $2^k$.
        \item \textit{Арифметический сдвиг вправо} (>> для signed): слева копируется знаковый бит.
        \item \textit{Циклический сдвиг}: вышедшие биты вставляются с противоположной стороны.
    \end{itemize}
\end{itemize}

Примеры (8 бит):
\begin{verbatim}
x = 0b01101011
x & 0b00001111 = 0b00001011  (выделение младших 4 бит)
x | 0b10000000 = 0b11101011  (установка старшего бита)
x ^ x = 0
x << 2 = 0b10101100
\end{verbatim}

\subsection{Двоичная арифметика с учётом ограниченной разрядности}
Сложение с переносами. Переполнение (overflow) при знаковых операциях фиксируется флагом OF. Пример: $127 + 1 = -128$ в 8-битном доп. коде.

Вычитание через доп. код: $A - B = A + (\sim B + 1)$.

Умножение и деление реализуются сдвигами и сложением/вычитанием.

Флаги состояния: ZF (ноль), SF (знак), CF (беззнаковое переполнение), OF (знаковое переполнение).

\subsection{Формат целых чисел}
В ЭВМ используются форматы с фиксированной точкой (1, 2, 4 байта) и двоично-десятичный формат (редко). Диапазоны:
\begin{itemize}
    \item 1 байт без знака: 0..255, со знаком: -128..127.
    \item 2 байта: 0..65535, со знаком: -32768..32767.
    \item 4 байта: 0..4294967295, со знаком: -2147483648..2147483647.
\end{itemize}

\subsection{Сдвиговый код (код со смещением)}
Используется в основном для хранения порядка чисел с плавающей точкой. Число $X$ представляется как беззнаковое $X + \text{bias}$. Для $n$ бит bias обычно выбирают $2^{n-1}$ или $2^{n-1}-1$. Ноль представляется единственным образом. Арифметика требует коррекции результата.

\section{Семинар №2. Целые числа: знаковые и беззнаковые. Диапазоны. Прямой, обратный, дополнительный коды. Оптимизация арифметики отрицательных чисел с помощью доп. кода. Побитовые операции. Двоичная арифметика.}

\subsection*{Задача 1 (диапазоны)}
Сколько различных чисел можно представить 12-битным беззнаковым целым? А 12-битным знаковым в дополнительном коде? Укажите минимальное и максимальное значение для каждого случая.

\textbf{Решение:}
\begin{itemize}
    \item Беззнаковое 12 бит: $2^{12} = 4096$ чисел; диапазон $[0, 4095]$.
    \item Знаковое доп. код: $2^{12} = 4096$; диапазон $[-2048, 2047]$.
\end{itemize}

\subsection*{Задача 2 (прямой, обратный, дополнительный коды)}
Представьте число $-37$ в 8-битных прямом, обратном и дополнительном кодах. Также в сдвиговом коде со смещением 128.

\textbf{Решение:}
Модуль $37_{10} = 00100101_2$.
\begin{itemize}
    \item Прямой: $10100101_2$.
    \item Обратный: инверсия $+37$: $11011010_2$.
    \item Дополнительный: $11011010_2 + 1 = 11011011_2$.
    \item Сдвиговый (bias=128): код = $-37+128=91 = 01011011_2$.
\end{itemize}

\subsection*{Задача 3 (сложение в дополнительном коде)}
Выполните сложение 8-битных чисел в дополнительном коде, интерпретируйте результат:
\begin{enumerate}
    \item[a)] $45 + (-23)$
    \item[b)] $100 + 30$ (проверьте флаги CF и OF)
    \item[c)] $(-128) + (-1)$
\end{enumerate}

\textbf{Решение:}
\begin{enumerate}
    \item[a)] $45=00101101_2$, $-23=11101001_2$. Сумма $=00010110_2=22_{10}$.
    \item[b)] $100=01100100_2$, $30=00011110_2$ $\to$ сумма $=10000010_2$ (беззнак. 130, знак. -126). CF=0, OF=1.
    \item[c)] $-128=10000000_2$, $-1=11111111_2$ $\to$ сумма $=01111111_2=127$. OF=1, CF=1.
\end{enumerate}

\subsection*{Задача 4 (побитовые операции)}
Дано $x = 0x5A$, $y = 0x3C$. Вычислите:
\begin{enumerate}
    \item[a)] $x \& y$
    \item[b)] $x | y$
    \item[c)] $x \oplus y$
    \item[d)] $\sim x$ (8 бит)
    \item[e)] $x << 2$ (логический)
    \item[f)] арифметический сдвиг вправо для знакового $x = 0xA6$ на 2.
\end{enumerate}

\textbf{Решение:}
$x=01011010_2$, $y=00111100_2$.
\begin{enumerate}
    \item[a)] $00011000_2 = 0x18$.
    \item[b)] $01111110_2 = 0x7E$.
    \item[c)] $01100110_2 = 0x66$.
    \item[d)] $10100101_2 = 0xA5$.
    \item[e)] $01101000_2 = 0x68$.
    \item[f)] $0xA6=10100110_2$; арифм. сдвиг вправо на 2: $11101001_2 = 0xE9$.
\end{enumerate}

\subsection*{Задача 5 (флаги и интерпретации)}
8-битные числа (доп. код). Выполните вычитание $80 - 50$ через сложение с доп. кодом. Укажите флаги.

\textbf{Решение:}
$80=01010000_2$, $-50=11001110_2$. Сумма $=1\,00011110$ (9 бит) $\to$ $00011110_2=30$. CF=1, OF=0.

\subsection*{Задача 6 (арифметика в сдвиговом коде)}
Два числа $A=-20$ и $B=35$ представлены в 8-битном сдвиговом коде со смещением 128. Выполните сложение и объясните коррекцию.

\textbf{Решение:}
Коды: $A= -20+128=108=01101100_2$; $B=35+128=163=10100011_2$. Сложение кодов даёт $00001111_2=15$ (без коррекции). Истинная сумма $=15$. Чтобы получить код суммы, нужно вычесть bias: $108+163-128=143=10001111_2$, что соответствует $15$.


\chapter{Тема №3. Числа с плавающей точкой.}
\section{Лекция №3. Числа с плавающей точкой. Мантисса и экспонента. Нормальная и нормализованная формы. Стандарт IEEE 754. Машинная точность.}

\subsection{Представление вещественных чисел: нормальная форма}
Вещественное число можно записать как $X = \pm M \times B^E$, где $M$ — мантисса, $B$ — основание (обычно 2), $E$ — экспонента.

Для однозначности мантиссу приводят к \textbf{нормализованной форме}: $1 \le |M| < B$. В двоичной системе целая часть мантиссы всегда равна 1 (кроме нуля и денормализованных). В памяти хранятся: знак (1 бит), экспонента со смещением, дробная часть мантиссы (без ведущей единицы).

В ЭВМ используется формат с плавающей точкой: число представляется в виде $A = \pm M \cdot S^{\pm P}$. Мантисса и порядок кодируются в фиксированных разрядах. Обычно число предварительно нормализуется: $S^{-1} \le |M| < 1$. Тогда целая часть мантиссы всегда 0, и её можно не хранить, что увеличивает точность.

Порядок кодируется в виде \textbf{характеристики} $Z = P + D$, где $D$ — смещение, подбираемое так, чтобы минимальный порядок давал характеристику 0. Например, в шестнадцатеричной нормализации $Z_{16} = 40_{16} + P_{16}$.

Пример (32 разряда, длина мантиссы 24, характеристики 7, знак 1): число $5.5_{10} = 5.8_{16} \cdot 16^0$. Нормализуем: $0.58_{16} \cdot 16^1$. Тогда $Z_{16}=40+1=41$, мантисса $0.58$, знак 0. Код: \texttt{41 58 00 00}$_{16}$.

\subsection{Стандарт IEEE 754}
Определяет форматы чисел с плавающей точкой.

\textbf{Одинарная точность (float, 32 бита)}:
\[
\begin{array}{|c|c|c|}
\hline
\text{Знак (1)} & \text{Экспонента (8)} & \text{Мантисса (23)} \\
\hline
\end{array}
\]
Смещение bias = 127. Истинная экспонента $E = e - 127$, где $e$ — хранимое беззнаковое целое.

\begin{itemize}
    \item $0 < e < 255$: нормализованное число: $(-1)^S \times (1.\text{дробная}) \times 2^{e-127}$.
    \item $e = 0$, мантисса = 0: $\pm 0$.
    \item $e = 0$, мантисса $\neq 0$: денормализованные: $(-1)^S \times (0.\text{дробная}) \times 2^{-126}$.
    \item $e = 255$, мантисса = 0: $\pm\infty$.
    \item $e = 255$, мантисса $\neq 0$: NaN.
\end{itemize}

\textbf{Двойная точность (double, 64 бита)}:
Смещение bias = 1023, экспонента 11 бит, мантисса 52 бита.

Пример кодирования $-12.625$ в float:
\begin{enumerate}
    \item $12.625_{10} = 1100.101_2 = 1.100101_2 \times 2^3$.
    \item $E=3$, $e=3+127=130=10000010_2$.
    \item Мантисса: $10010100000000000000000$ (23 бита).
    \item Знак = 1.
    \item Код: \texttt{1 10000010 10010100000000000000000} = \texttt{0xC14A0000}.
\end{enumerate}

\subsection{Машинная точность (машинный эпсилон)}
$\varepsilon$ — разница между 1.0 и наименьшим представимым числом, большим 1.0.
Для float: $\varepsilon = 2^{-23} \approx 1.19 \times 10^{-7}$.
Для double: $\varepsilon = 2^{-52} \approx 2.22 \times 10^{-16}$.

Относительная погрешность любого представимого числа не превосходит $\varepsilon/2$ при округлении к ближайшему.

\subsection{Особенности арифметики с плавающей точкой}
\begin{itemize}
    \item Округление к ближайшему чётному.
    \item Поглощение: при сложении чисел с сильно разными порядками меньшее может быть потеряно.
    \item Неассоциативность: $(a+b)+c \neq a+(b+c)$.
    \item Сравнение на равенство не рекомендуется.
\end{itemize}

\section{Семинар №3. Числа с плавающей точкой.}

\subsection*{Задача 1 (Нормализация и представление)}
Число $23.375_{10}$ представьте в нормализованной двоичной форме, а затем в формате float IEEE 754. Укажите шестнадцатеричное представление.

\textbf{Решение:}
$23.375 = 10111.011_2 = 1.0111011_2 \times 2^4$. $E=4$, $e=4+127=131=10000011_2$. Мантисса $0111011\ldots$ (23 бита). Код: \texttt{0 10000011 01110110000000000000000} = \texttt{0x41BB0000}.

\subsection*{Задача 2 (Декодирование float)}
Дано 32-битное число в hex: \texttt{0xC0A00000}. Определите десятичное число.

\textbf{Решение:}
\texttt{C0A00000} = $1100\,0000\,1010\,0000\ldots_2$. Знак=1, экспонента $e=129$, $E=2$, мантисса $=1.01_2=1.25$. Число = $-1.25 \times 2^2 = -5.0$.

\subsection*{Задача 3 (Машинный эпсилон)}
Объясните, почему $0.1$ и $0.2$ не представимы точно в float. Чему равен $\varepsilon$ для float?

\textbf{Решение:}
$0.1_{10}$ в двоичной — периодическая дробь. $\varepsilon = 2^{-23} \approx 1.1920929\times10^{-7}$.

\subsection*{Задача 4 (Особые значения)}
Что будет результатом: $1.0/0.0$, $0.0/0.0$, $\infty-\infty$, $10^{38}\cdot10^{38}$ (float)? Определите битовое представление бесконечности и NaN.

\textbf{Решение:}
$+\infty$ (e=255,m=0), NaN (e=255,m$\neq$0), NaN, $\infty$. $+\infty$: \texttt{0 11111111 000...0}. NaN: \texttt{0 11111111 100...0}.

\subsection*{Задача 5 (Поглощение и неассоциативность)}
Покажите на примере float, что $(1.0 + 10^{-8}) + 10^{-8}$ может отличаться от $1.0 + (10^{-8} + 10^{-8})$.

\textbf{Решение:}
$\varepsilon \approx 1.19\times10^{-7}$. $10^{-8}$ меньше $\varepsilon$, поэтому при сложении с 1.0 она теряется. Оба выражения дадут 1.0. Для демонстрации лучше взять $10^{-7}$ (чуть меньше $\varepsilon$). Тогда первое выражение даст 1.0, а второе — $1.0 + 2\cdot10^{-7}$ может дать $1.00000012$ (если $2\cdot10^{-7} > \varepsilon/2$). Неассоциативность проявляется.


\chapter{Тема №4. Логические основы функционирования ЭВМ.}
\section{Лекция №4. Логические операции, законы, графическое представление, поразрядные операции, сдвиги.}

\subsection{Логические операции и функции}
Алгебра логики (булева алгебра) оперирует логическими функциями, аргументы и значения которых принимают только два значения: «истина» (1) или «ложь» (0).

Основные логические операции:
\begin{itemize}
    \item \textbf{Инверсия (НЕ, NOT)}: $!x$ — результат противоположен операнду.
    \item \textbf{Конъюнкция (И, AND)}: $x_1 \land x_2$ — истинна, только если оба операнда истинны.
    \item \textbf{Дизъюнкция (ИЛИ, OR)}: $x_1 \lor x_2$ — истинна, если хотя бы один операнд истинен.
\end{itemize}

Таблицы истинности:

\begin{center}
\begin{tabular}{c|c}
$x$ & $\neg x$ \\
\hline
0 & 1 \\
1 & 0
\end{tabular}
\qquad
\begin{tabular}{cc|c}
$x_1$ & $x_2$ & $x_1 \land x_2$ \\
\hline
0 & 0 & 0 \\
0 & 1 & 0 \\
1 & 0 & 0 \\
1 & 1 & 1
\end{tabular}
\qquad
\begin{tabular}{cc|c}
$x_1$ & $x_2$ & $x_1 \lor x_2$ \\
\hline
0 & 0 & 0 \\
0 & 1 & 1 \\
1 & 0 & 1 \\
1 & 1 & 1
\end{tabular}
\end{center}

Дополнительные операции:
\begin{itemize}
    \item \textbf{Исключающее ИЛИ (XOR)}: $x \oplus y$ — истинна, если операнды различны.
    \item \textbf{Эквивалентность (XNOR)}: $x = y$ — истинна, если операнды равны.
    \item \textbf{Импликация}: $x \rightarrow y$ — ложна, только если $x=1, y=0$.
\end{itemize}

Порядок выполнения логических операций: сначала отрицание, затем конъюнкция, затем дизъюнкция. Скобки изменяют порядок.

\subsection{Законы алгебры логики}
\begin{itemize}
    \item Коммутативность: $a\lor b = b\lor a$, $a\land b = b\land a$.
    \item Ассоциативность: $(a\lor b)\lor c = a\lor(b\lor c)$, аналогично для $\land$.
    \item Дистрибутивность: $a\land(b\lor c) = (a\land b)\lor(a\land c)$; $a\lor(b\land c) = (a\lor b)\land(a\lor c)$.
    \item Законы де Моргана: $\overline{a\land b} = \overline{a}\lor\overline{b}$; $\overline{a\lor b} = \overline{a}\land\overline{b}$.
    \item Двойное отрицание: $\overline{\overline{a}} = a$.
    \item Закон исключения: $a\lor\overline{a}=1$.
    \item Закон противоречия: $a\land\overline{a}=0$.
    \item Свойства: $0\land a=0$, $1\land a=a$, $0\lor a=a$, $1\lor a=1$, $a\land a\land\ldots=a$, $a\lor a\lor\ldots=a$.
\end{itemize}

Следствия: склеивание, неполное склеивание, поглощение.

Минимизация логических выражений сокращает количество операций и элементов в схеме.

\subsection{Графическое представление логических функций}
В технике логические элементы обозначаются стандартными значками (И, ИЛИ, НЕ, И-НЕ, ИЛИ-НЕ и др.). Например, элемент И-НЕ (штрих Шеффера) и ИЛИ-НЕ (стрелка Пирса) являются функционально полными базисами.

Пример реализации функции $Y = \overline{x_1 + x_1 x_2 x_3}$ с использованием инвертора и элементов И, ИЛИ.

\subsection{Поразрядные логические операции и операции сдвигов}
Поразрядные операции применяются к каждому биту кода независимо. Они используются для работы с битовыми полями (маски).

Типовые операции:
\begin{itemize}
    \item Инверсия битов (NOT).
    \item Битовая конъюнкция (AND) — для обнуления или выделения битов.
    \item Битовая дизъюнкция (OR) — для установки битов в 1.
    \item Битовая неравнозначность (XOR) — для инверсии выбранных битов.
\end{itemize}

Маска — специальный код, в котором единицы стоят в тех битах, которые нужно изменить или проверить.

Примеры:
\begin{enumerate}
    \item Проверка состояния битов 1 и 5 (нумерация справа налево с 1): маска $10001_2$, AND с исходным кодом.
    \item Установка битов 1 и 2 в 0: маска с нулями в этих битах, AND.
    \item Установка битов 1 и 2 в 1: маска с единицами, OR.
    \item Инверсия битов 1,2,5: маска с единицами, XOR.
    \item Занесение кода в поле: обнуление поля маской AND, затем OR с кодом.
\end{enumerate}

Сдвиговые операции:
\begin{itemize}
    \item \textbf{Логический сдвиг:} освободившиеся биты заполняются нулями. Сдвиг влево на $n$ равносилен умножению на $2^n$ (беззнаковое), вправо — делению на $2^n$.
    \item \textbf{Арифметический сдвиг:} при сдвиге вправо освободившиеся биты заполняются значением знакового бита (сохраняет знак).
    \item \textbf{Циклический сдвиг:} вытолкнутые биты вставляются с противоположной стороны.
\end{itemize}

\section{Семинар №4. Логические основы.}

\subsection*{Задача 1 (логические функции)}
Какие сигналы будут на выходах S, Q, R при заданных входных сигналах комбинационной схемы? (схема не приведена, но можно описать).

\subsection*{Задача 2 (преобразование функций)}
Преобразовать функцию к виду, используя законы алгебры логики.

\subsection*{Задача 3 (таблица истинности)}
Какой аналитической записи соответствует таблица истинности? (таблица приведена в Word).

\subsection*{Задача 4 (поразрядные операции)}
Даны два числа. Выполнить поразрядные операции AND, OR, XOR, NOT.


\chapter{Тема №5. Кодирование команд и символьных данных.}
\section{Лекция №5. Кодирование команд. Кодирование символьных данных.}

\subsection{Кодирование команд}
Машинная команда содержит код операции и адресную часть. Код операции определяет действие, адреса указывают операнды и результат.

Длина кода операции $p$ бит позволяет закодировать до $2^p$ различных команд. Оптимальная длина: $p = \lceil \log_2 P \rceil$, где $P$ — число команд.

Пример трехадресной команды: \texttt{КОП А1 А2 А3}.

\subsection{Кодирование символьных данных}
Символы кодируются целыми беззнаковыми числами. Используются кодовые таблицы:

\begin{itemize}
    \item \textbf{ASCII} (7 или 8 бит): первые 128 символов — стандартные (латиница, цифры, знаки), расширение (128-255) — для национальных алфавитов.
    \item \textbf{ISO 8859} — серия 8-битных кодировок, включая кириллицу (ISO 8859-5).
    \item \textbf{Unicode} (UTF-8, UTF-16, UTF-32) — универсальная кодировка.
    \begin{itemize}
        \item UTF-8: переменная длина (1-4 байта), совместима с ASCII для кодов 0-127.
        \item UTF-16: 2 или 4 байта.
        \item UTF-32: 4 байта фиксированно.
    \end{itemize}
    \item Кириллические кодировки: Windows-1251, CP-866, KOI-8R и др.
\end{itemize}

\section{Семинар №5. Кодирование команд и символьных данных.}
(Задач по этой теме в источнике нет, поэтому семинар пустой.)


\chapter{Задачи для самостоятельного решения}
\section*{Условия задач (без решений)}

Ниже приведены все задачи из курса. Решения приведены в соответствующих семинарах.

\subsection*{Задачи из семинаров}
\paragraph{Семинар №1}
\begin{enumerate}
    \item Переведите число $250_{10}$ в двоичную, восьмеричную и шестнадцатеричную системы.
    \item Переведите числа $1100101_2$, $145_8$, $A9_{16}$ в десятичную систему.
    \item Переведите двоичное число $110111001010_2$ в восьмеричную и шестнадцатеричную. Переведите $2F5_{16}$ в двоичную и восьмеричную.
    \item Выполните действия:
    \begin{enumerate}
        \item $1011_2 + 1101_2$
        \item $11001_2 - 1011_2$
        \item $23_8 \times 5_8$
        \item $1A3_{16} + 4F_{16}$
        \item $200_8 / 7_8$ (целочисленное деление)
    \end{enumerate}
\end{enumerate}

\paragraph{Семинар №2}
\begin{enumerate}
    \setcounter{enumi}{5}
    \item Сколько различных чисел можно представить 12-битным беззнаковым целым? А 12-битным знаковым в дополнительном коде? Укажите диапазоны.
    \item Представьте число $-37$ в 8-битных прямом, обратном, дополнительном кодах и в сдвиговом коде со смещением 128.
    \item Выполните сложение 8-битных чисел в дополнительном коде:
    \begin{enumerate}
        \item $45 + (-23)$
        \item $100 + 30$ (укажите флаги CF и OF)
        \item $(-128) + (-1)$
    \end{enumerate}
    \item Дано $x = 0x5A$, $y = 0x3C$. Вычислите:
    \begin{enumerate}
        \item $x \& y$
        \item $x | y$
        \item $x \oplus y$
        \item $\sim x$ (8 бит)
        \item $x << 2$ (логический сдвиг)
        \item арифметический сдвиг вправо для знакового $x = 0xA6$ на 2.
    \end{enumerate}
    \item Выполните вычитание $80 - 50$ через сложение с дополнительным кодом. Укажите флаги.
    \item Два числа $A=-20$ и $B=35$ представлены в 8-битном сдвиговом коде со смещением 128. Выполните сложение и объясните коррекцию.
\end{enumerate}

\paragraph{Семинар №3}
\begin{enumerate}
    \setcounter{enumi}{12}
    \item Число $23.375_{10}$ представьте в нормализованной двоичной форме и в формате float IEEE 754. Укажите шестнадцатеричное представление.
    \item Дано 32-битное число в hex: \texttt{0xC0A00000}. Определите десятичное число.
    \item Объясните, почему $0.1$ и $0.2$ не представимы точно в float. Чему равен $\varepsilon$ для float?
    \item Что будет результатом: $1.0/0.0$, $0.0/0.0$, $\infty-\infty$, $10^{38}\cdot10^{38}$ (float)? Определите битовое представление бесконечности и NaN.
    \item Покажите на примере float неассоциативность сложения.
\end{enumerate}

\paragraph{Семинар №4 (логика)}
\begin{enumerate}
    \setcounter{enumi}{17}
    \item Какие сигналы будут на выходах S, Q, R при заданных входных сигналах комбинационной схемы? (задача из Word №4)
    \item Преобразовать функцию $Y = \overline{x_1} \vee \overline{x_2} \vee \overline{x_3}$ к виду, используя законы алгебры логики.
    \item Какой аналитической записи логической функции соответствует таблица истинности? (задача из Word №7)
\end{enumerate}

\subsection*{Задачи из Word-документа (дополнительные)}
\begin{enumerate}
    \setcounter{enumi}{20}
    \item Перевести $A_{16}=852.8$ в $A_{10}$ и далее в $A_8$.
    \item Выполнить перевод $A_{16}=852.8$ в $A_8$ и далее в $A_{10}$. Сравните результат с задачей 1.
    \item Вычислить сумму и разность двух чисел в СС$_5$: $123_5$ и $12_5$. Перевести исходные числа, сумму и разность в десятичную систему. Проверить правильность вычислений.
    \item (задача на логику, дублируется с №18)
    \item Преобразовать функцию $F = x_1 x_2 \vee \overline{x_1} x_3$ к виду $F = (x_1 \vee x_3)(x_2 \vee \overline{x_1})$.
    \item Преобразовать функцию $F = (x_1 \vee x_2)(x_1 \vee \overline{x_2})$ к виду $F = x_1$.
    \item (задача на таблицу истинности, дублируется с №19)
    \item Даны два числа: $A=-30$ (десятичное) и $B=-20$ (шестнадцатеричное). Представлены в дополнительном коде в 16-разрядном формате. Выполнить $A+B$ и определить восьмеричный дополнительный код результата.
    \item Даны два числа: $A=-20$ (десятичное) и $B=-11$ (шестнадцатеричное). Выполнить $A-B$ в дополнительном коде (16 разрядов) и определить восьмеричный дополнительный код результата.
    \item Представить дополнительный код числа $A_{10}=-423$ в СС$_8$ при $N=16$.
    \item Цветное растровое изображение с палитрой 65536 цветов имеет размер 100×100 точек. Какой объем видеопамяти (в Кбайтах) занимает это изображение?
\end{enumerate}

\end{document}
